\documentclass[UTF8]{ctexart}

\usepackage[a4paper,margin=1in]{geometry}
\usepackage{amsmath,amssymb,mathtools,bm}
\usepackage{booktabs,longtable,array}
\usepackage{hyperref}
\usepackage{enumitem}

\hypersetup{
    colorlinks=true,
    linkcolor=black,
    urlcolor=blue,
    citecolor=black
}

\setlist[itemize]{leftmargin=2em}
\setlist[enumerate]{leftmargin=2em}

\title{\texttt{mra.py} 的数学原理与数据流说明}
\author{}
\date{\today}

\begin{document}

\maketitle
\tableofcontents
\newpage

\section{文档范围与实现基准}

本文档以仓库中的 \texttt{mra.py} 为唯一实现基准，目标是完整说明该脚本的：
\begin{itemize}
    \item 数学建模思路；
    \item 各模块的张量形状变化；
    \item 训练、推理与异常评分的完整数据流；
    \item 与常见论文表述不同但在代码中真实存在的实现细节。
\end{itemize}

换言之，本文档描述的是\textbf{当前代码的真实行为}，而不仅仅是一个理想化的模型设计图。

\section{符号、张量形状与超参数}

\subsection{核心符号}

\begin{longtable}{>{\raggedright\arraybackslash}p{0.17\linewidth} >{\raggedright\arraybackslash}p{0.23\linewidth} >{\raggedright\arraybackslash}p{0.50\linewidth}}
\toprule
符号 & 形状 & 含义 \\
\midrule
\endhead
$B$ & 标量 & 批大小（batch size） \\
$S$ & 标量 & 窗口长度，代码中默认 $S=60$ \\
$N$ & 标量 & 变量个数，即 CSV 列数，运行时由数据决定 \\
$\bm{X}$ & $\mathbb{R}^{B \times S \times N}$ & 一个批次的窗口化输入序列 \\
$\bm{M}$ & $\{0,1\}^{B \times S \times N}$ & 缺失掩码，\textbf{$1$ 表示 missing，$0$ 表示 observed} \\
$\bm{A}$ & $\mathbb{R}^{N \times N}$ & GraphLearner 生成的邻接矩阵 \\
$\bm{H}^{\text{gcn}}$ & $\mathbb{R}^{B \times S \times N}$ & 图卷积分支输出 \\
$\bm{H}^{\text{tcn}}$ & $\mathbb{R}^{B \times S \times N}$ & 多尺度时序卷积分支输出 \\
$\bm{H}^{\text{time}}$ & $\mathbb{R}^{B \times S \times N}$ & 时域分支综合表示 \\
$\bm{H}^{\text{freq}}$ & $\mathbb{R}^{B \times S \times N}$ & 频域分支输出 \\
$\bm{X}^{\text{imp}}$ & $\mathbb{R}^{B \times S \times N}$ & 门控融合后的插补表示 \\
$\bm{Z}$ & $\mathbb{R}^{B \times S \times d}$ & Transformer 隐空间表示，代码中 $d=64$ \\
$\hat{\bm{X}}$ & $\mathbb{R}^{B \times S \times N}$ & 最终重构输出 \\
\bottomrule
\end{longtable}

\subsection{来自代码的固定超参数}

\begin{longtable}{>{\raggedright\arraybackslash}p{0.35\linewidth} >{\raggedright\arraybackslash}p{0.15\linewidth} >{\raggedright\arraybackslash}p{0.35\linewidth}}
\toprule
项目 & 数值 & 说明 \\
\midrule
\endhead
窗口长度 $S$ & $60$ & \texttt{SEQ\_LEN = 60} \\
滑窗步长 & $1$ & \texttt{DatasetBuilder(stride=1)} \\
图嵌入维度 & $16$ & \texttt{GraphLearner(embed\_dim=16)} \\
$\alpha$ & $3.0$ & 图嵌入缩放因子 \\
TCN 卷积核 & $\{3,5,9\}$ & 主模型中实际使用的三种尺度 \\
TCN 模式 & non-causal & 主模型中 \texttt{causal=False} \\
Transformer 隐维 $d$ & $64$ & \texttt{d\_model=64} \\
注意力头数 & $4$ & \texttt{nhead=4} \\
Transformer 层数 & $2$ & \texttt{num\_layers=2} \\
FFN 隐维 & $128$ & \texttt{dim\_feedforward=128} \\
Dropout & $0.1$ & Transformer 编码层内置 dropout \\
优化器学习率 & $10^{-3}$ & Adam 初始学习率 \\
调度器 factor & $0.5$ & \texttt{ReduceLROnPlateau} 衰减倍数 \\
调度器 patience & $10$ & 10 个 epoch 不改进则降学习率 \\
训练轮数 & $3$ & 代码中当前写死为 \texttt{for epoch in range(3)} \\
批大小 & $32$ & 训练和推理默认均为 32 \\
频域损失权重 & $0.1$ & \texttt{freq\_weight} 默认值 \\
图熵正则权重 & $0.1$ & \texttt{sparsity\_weight} 默认值 \\
\bottomrule
\end{longtable}

\section{整体流程总览}

\subsection{训练期总流程}

训练时，一个窗口样本从数据到损失的大流程可以写成：
\begin{align}
    (\bm{X}, \bm{M}) 
    &\xrightarrow{\text{随机再遮挡}}
    (\bm{X}^{\text{in}}, \bm{M}^{\text{in}}, \bm{T}) \\
    &\xrightarrow{\text{AGF-ADNet}}
    (\hat{\bm{X}}, \bm{A}, \bm{X}^{\text{imp}}) \\
    &\xrightarrow{\text{dual\_domain\_loss}}
    \mathcal{L}_{\text{rec}} + 0.1\,\mathcal{L}_{\text{freq}} + 0.1\,\mathcal{L}_{\text{ent}}
\end{align}
其中：
\begin{itemize}
    \item $\bm{X}$ 是原始窗口；
    \item $\bm{M}$ 是原始缺失掩码；
    \item $\bm{M}^{\text{in}}$ 是加入随机遮挡后的输入掩码；
    \item $\bm{T}$ 是本批次真正参与监督的目标位置掩码；
    \item $\bm{A}$ 是学习到的图结构；
    \item $\hat{\bm{X}}$ 是最终重构输出。
\end{itemize}

\subsection{推理期总流程}

推理时没有随机再遮挡，因此流程更简单：
\begin{align}
    (\bm{X}, \bm{M})
    \xrightarrow{\text{apply\_missing\_mask}}
    \bm{X}^{\text{in}}
    \xrightarrow{\text{AGF-ADNet}}
    \hat{\bm{X}}
    \xrightarrow{\text{masked MSE}}
    \text{Score}
\end{align}

测试阈值并不是通过验证集调参得到，而是直接由训练集分数的均值与标准差估计：
\begin{equation}
    \tau = \mu_{\text{train}} + \sigma_{\text{train}}.
\end{equation}

\section{数据准备与窗口构造}

\subsection{目录读取与掩码定义}

\texttt{DatasetBuilder.load\_dir()} 从指定目录下读取一组 CSV 文件，并按文件名排序后顺序拼接。设单个原始时间序列为
\begin{equation}
    \bm{D} \in \mathbb{R}^{T \times N},
\end{equation}
其中 $T$ 为时间长度，$N$ 为变量数。

代码对缺失值使用如下掩码约定：
\begin{equation}
    M_{t,n} =
    \begin{cases}
        1, & D_{t,n}\text{ 为缺失值} \\
        0, & D_{t,n}\text{ 为观测值}
    \end{cases}
\end{equation}
这个定义与很多论文中“$1$ 表示 observed”正好相反，因此后续阅读代码时必须持续记住该语义。

\subsection{标准化}

标准化前，代码先把缺失值用 $0$ 临时填充：
\begin{equation}
    \bar{D}_{t,n} =
    \begin{cases}
        0, & D_{t,n}\text{ 缺失} \\
        D_{t,n}, & \text{否则}
    \end{cases}
\end{equation}

然后仅在训练集上拟合 \texttt{StandardScaler}：
\begin{align}
    \mu_n &= \frac{1}{T_{\text{train}}}\sum_{t=1}^{T_{\text{train}}}\bar{D}_{t,n}, \\
    \sigma_n^2 &= \frac{1}{T_{\text{train}}}\sum_{t=1}^{T_{\text{train}}}(\bar{D}_{t,n}-\mu_n)^2,
\end{align}
并对训练集和测试集统一变换：
\begin{equation}
    D^{\text{scaled}}_{t,n} = \frac{\bar{D}_{t,n} - \mu_n}{\sigma_n}.
\end{equation}

\paragraph{实现含义}
这里的标准化对象不是“带 NaN 的真实数据”，而是“先用 0 填过 NaN 的数组”。因此，缺失位置在标准化之后仍然会带有一个数值占位符；真正让缺失位置退出模型输入的是后面的 \texttt{apply\_missing\_mask()}。

\subsection{滑动窗口}

设标准化后的序列记为 $\bm{D}^{\text{scaled}} \in \mathbb{R}^{T \times N}$，窗口长度为 $S$，步长为 $\delta=1$。代码以时间索引 $i$ 为窗口终点，构造第 $i$ 个窗口：
\begin{equation}
    \bm{X}^{(i)} =
    \begin{cases}
        [\underbrace{\bm{d}_0,\dots,\bm{d}_0}_{S-i-1},\bm{d}_0,\bm{d}_1,\dots,\bm{d}_i], & i < S \\
        [\bm{d}_{i-S+1}, \dots, \bm{d}_i], & i \ge S
    \end{cases}
\end{equation}
其中 $\bm{d}_t \in \mathbb{R}^{N}$ 为第 $t$ 个时间步的向量。

掩码窗口以完全相同的方式构造：
\begin{equation}
    \bm{M}^{(i)} =
    \begin{cases}
        [\underbrace{\bm{m}_0,\dots,\bm{m}_0}_{S-i-1},\bm{m}_0,\bm{m}_1,\dots,\bm{m}_i], & i < S \\
        [\bm{m}_{i-S+1}, \dots, \bm{m}_i], & i \ge S
    \end{cases}
\end{equation}

最终一个批次输入模型的张量形状就是：
\begin{equation}
    \bm{X}, \bm{M} \in \mathbb{R}^{B \times S \times N}.
\end{equation}

\section{模型主体：AGF-ADNet}

\subsection{前向传播总式}

对于一个输入窗口批次，\texttt{AGF\_ADNet.forward(x, mask)} 的主流程可以抽象为：
\begin{align}
    \bm{A} &= \operatorname{GraphLearner}() \\
    \bm{H}^{\text{time}} &= \operatorname{TimeBranch}(\bm{X}^{\text{in}}, \bm{A}) \\
    \bm{H}^{\text{freq}} &= \operatorname{FreqBranch}(\bm{X}^{\text{in}}) \\
    \bm{X}^{\text{imp}} &= \operatorname{GatedFusion}(\bm{H}^{\text{time}}, \bm{H}^{\text{freq}}) \\
    \hat{\bm{X}} &= \operatorname{TransformerHead}(\bm{X}^{\text{filled}})
\end{align}
其中 $\bm{X}^{\text{filled}}$ 的\textbf{理想形式}应当是“观测位置保留原值，缺失位置使用插补值”，但当前代码的真实行为与理想公式存在差异，后文会专门说明。

\subsection{图结构学习模块：GraphLearner}

\subsubsection{参数化}

GraphLearner 使用两组可学习节点嵌入：
\begin{equation}
    \bm{E}_1, \bm{E}_2 \in \mathbb{R}^{N \times d_e}, \qquad d_e = 16.
\end{equation}

它们先经过缩放与双曲正切非线性：
\begin{align}
    \bm{M}_1 &= \tanh(\alpha \bm{E}_1), \\
    \bm{M}_2 &= \tanh(\alpha \bm{E}_2),
\end{align}
其中 $\alpha = 3.0$。

\subsubsection{邻接矩阵计算}

原始邻接分数为：
\begin{equation}
    \bm{A}^{\text{raw}} = \operatorname{ReLU}(\bm{M}_1 \bm{M}_2^\top).
\end{equation}

随后进行逐行归一化：
\begin{equation}
    A_{ij} = \frac{A^{\text{raw}}_{ij}}{\sum_{k=1}^{N}A^{\text{raw}}_{ik} + \varepsilon},
    \qquad \varepsilon = 10^{-8}.
\end{equation}

\paragraph{为什么不是 softmax}
代码刻意避免使用 softmax，而采用 ReLU 后再做行归一化。原因是 softmax 会让所有边权都严格大于 0，图结构难以变得真正尖锐；当前实现后面还会对邻接矩阵施加熵约束，因此需要允许某些边直接变为 0。

\subsection{时域分支}

\subsubsection{GCNLayer}

输入先扩展一个特征维：
\begin{equation}
    \bm{X}^{\text{in}} \in \mathbb{R}^{B \times S \times N}
    \;\longrightarrow\;
    \bm{X}^{\text{in,4D}} \in \mathbb{R}^{B \times S \times N \times 1}.
\end{equation}

图卷积层先做线性映射：
\begin{equation}
    \tilde{\bm{X}}_{b,s,n,:} = \bm{W}\bm{X}^{\text{in,4D}}_{b,s,n,:} + \bm{b},
\end{equation}
然后按邻接矩阵在节点维聚合：
\begin{equation}
    H^{\text{gcn}}_{b,s,n,f}
    =
    \sum_{m=1}^{N} A_{n,m}\,\tilde{X}_{b,s,m,f}.
\end{equation}

在代码中该操作由
\begin{equation}
    \texttt{torch.einsum("nm,bsmd->bsnd", adj, x)}
\end{equation}
实现。

线性层的输入输出维度都是 $1$，因此最终再 squeeze 回到
\begin{equation}
    \bm{H}^{\text{gcn}} \in \mathbb{R}^{B \times S \times N},
\end{equation}
并执行 LayerNorm：
\begin{equation}
    \bm{H}^{\text{gcn}} \leftarrow \operatorname{LayerNorm}(\bm{H}^{\text{gcn}}).
\end{equation}

\subsubsection{MultiScaleTCN}

输入先转置到通道优先格式：
\begin{equation}
    \bm{X}^{\text{in}} \in \mathbb{R}^{B \times S \times N}
    \;\longrightarrow\;
    \bm{X}^{\text{tcn}} \in \mathbb{R}^{B \times N \times S}.
\end{equation}

对每个卷积核尺寸 $k \in \{3,5,9\}$，代码使用一层 depthwise Conv1D：
\begin{equation}
    \bm{Y}^{(k)} = \operatorname{Conv1D}_{k,\;groups=N}(\operatorname{Pad}(\bm{X}^{\text{tcn}})).
\end{equation}

因为在主模型中 \texttt{causal=False}，所以 padding 采用对称方式。对于奇数卷积核，有：
\begin{equation}
    p_{\text{left}} = \frac{k-1}{2}, \qquad
    p_{\text{right}} = \frac{k-1}{2}.
\end{equation}

每个节点都是独立卷积，不发生跨节点通道混合，因此第 $n$ 个节点、第 $s$ 个时间步的输出可写为：
\begin{equation}
    Y^{(k)}_{b,n,s}
    =
    \sum_{i=0}^{k-1} w^{(k)}_{n,i}\,\tilde{X}^{\text{tcn}}_{b,n,s+i-p_{\text{left}}}
    + b^{(k)}_n.
\end{equation}

三种尺度输出堆叠为：
\begin{equation}
    \bm{Y}^{\text{stack}} \in \mathbb{R}^{B \times N \times S \times 3},
\end{equation}
再经过一个线性层 $3 \rightarrow 1$ 融合：
\begin{equation}
    \bm{H}^{\text{tcn}} = \operatorname{Linear}_{3\rightarrow 1}(\bm{Y}^{\text{stack}}).
\end{equation}

最后转回 $(B,S,N)$，并与 GCN 分支相加：
\begin{equation}
    \bm{H}^{\text{time}} = \bm{H}^{\text{gcn}} + \bm{H}^{\text{tcn}}.
\end{equation}

\subsection{频域分支：FrequencyImputer}

\subsubsection{时域到频域}

输入先转为 $(B,N,S)$：
\begin{equation}
    \bm{X}^{\text{freq-in}} = \operatorname{permute}(\bm{X}^{\text{in}}) \in \mathbb{R}^{B \times N \times S}.
\end{equation}

然后沿时间维做实数快速傅里叶变换：
\begin{equation}
    \bm{X}_f = \operatorname{rFFT}(\bm{X}^{\text{freq-in}}) \in \mathbb{C}^{B \times N \times F},
    \qquad F = \left\lfloor \frac{S}{2} \right\rfloor + 1.
\end{equation}

代码中还计算了：
\begin{align}
    \bm{R} &= \Re(\bm{X}_f), \\
    \bm{I} &= \Im(\bm{X}_f), \\
    |\bm{X}_f| &= \text{magnitude}, \\
    \angle \bm{X}_f &= \text{phase}.
\end{align}

\paragraph{实现注记}
\texttt{magnitude} 与 \texttt{phase} 在当前代码中被计算出来，但后续并未参与损失或前向增强过程。真正进入神经网络的是实部与虚部的拼接特征。

\subsubsection{频率注意力与增强}

构造频域实值特征：
\begin{equation}
    \bm{Z}_f = [\bm{R};\bm{I}] \in \mathbb{R}^{B \times N \times 2F}.
\end{equation}

注意力网络输出频率权重：
\begin{equation}
    \bm{\alpha}
    =
    \sigma\!\left(
        \bm{W}^{\text{att}}_2
        \,\phi\!\left(
            \bm{W}^{\text{att}}_1 \bm{Z}_f + \bm{b}^{\text{att}}_1
        \right)
        + \bm{b}^{\text{att}}_2
    \right)
    \in [0,1]^{B \times N \times F},
\end{equation}
其中 $\phi(\cdot)$ 表示 ReLU。

增强网络输出：
\begin{equation}
    \hat{\bm{Z}}_f
    =
    \bm{W}^{\text{enh}}_2
    \,\phi\!\left(
        \bm{W}^{\text{enh}}_1 \bm{Z}_f + \bm{b}^{\text{enh}}_1
    \right)
    + \bm{b}^{\text{enh}}_2
    \in \mathbb{R}^{B \times N \times 2F}.
\end{equation}

将其拆成实部增强量与虚部增强量：
\begin{align}
    \hat{\bm{R}} &= \hat{\bm{Z}}_f[:, :, 1:F], \\
    \hat{\bm{I}} &= \hat{\bm{Z}}_f[:, :, F+1:2F].
\end{align}

更准确地说，若按 0-based 编号，代码切片实际是：
\begin{align}
    \hat{\bm{R}} &= \hat{\bm{Z}}_f[:, :, 0:F], \\
    \hat{\bm{I}} &= \hat{\bm{Z}}_f[:, :, F:2F].
\end{align}

频率注意力逐元素作用于增强量：
\begin{align}
    \bm{R}^{\text{att}} &= \hat{\bm{R}} \odot \bm{\alpha}, \\
    \bm{I}^{\text{att}} &= \hat{\bm{I}} \odot \bm{\alpha}.
\end{align}

随后以残差方式修正原始频谱：
\begin{equation}
    \bm{X}^{\text{enh}}_f = \bm{X}_f + \left(\bm{R}^{\text{att}} + j\,\bm{I}^{\text{att}}\right).
\end{equation}

\subsubsection{频域回到时域}

最后做逆实傅里叶变换并转回原排列：
\begin{align}
    \bm{H}^{\text{freq,raw}} &= \operatorname{irFFT}(\bm{X}^{\text{enh}}_f) \in \mathbb{R}^{B \times N \times S}, \\
    \bm{H}^{\text{freq}} &= \operatorname{LayerNorm}(\operatorname{permute}^{-1}(\bm{H}^{\text{freq,raw}})) \in \mathbb{R}^{B \times S \times N}.
\end{align}

\subsection{门控融合：GatedFusion}

将时域和频域分支分别转到 $(B,N,S)$，然后拼接：
\begin{equation}
    \bm{C} = [\bm{H}^{\text{time}};\bm{H}^{\text{freq}}] \in \mathbb{R}^{B \times 2N \times S}.
\end{equation}

门控网络结构是：
\begin{equation}
    \bm{Z}_g
    =
    \sigma\!\Big(
        \operatorname{Conv1D}_{1\times 1}
        \big(
            \phi(\operatorname{Conv1D}_{k=3}(\bm{C}))
        \big)
    \Big)
    \in [0,1]^{B \times N \times S}.
\end{equation}

转回 $(B,S,N)$ 后，逐元素融合：
\begin{equation}
    \bm{X}^{\text{imp}}
    =
    \operatorname{LayerNorm}
    \big(
        \bm{Z}_g \odot \bm{H}^{\text{time}}
        +
        (1-\bm{Z}_g)\odot \bm{H}^{\text{freq}}
    \big).
\end{equation}

\paragraph{解释}
当某个位置的门值接近 $1$ 时，模型更相信时域分支；当门值接近 $0$ 时，更相信频域分支。

\subsection{Transformer 重构头}

\subsubsection{输入投影}

理论上，融合后的插补序列应形成 $\bm{X}^{\text{filled}}$，再送入 Transformer。代码中的投影为：
\begin{equation}
    \bm{Z}_0 = \bm{X}^{\text{filled}}\bm{W}_{\text{in}} + \bm{b}_{\text{in}} + \bm{P},
\end{equation}
其中：
\begin{equation}
    \bm{W}_{\text{in}} \in \mathbb{R}^{N \times d}, \qquad
    \bm{P} \in \mathbb{R}^{1 \times S \times d}, \qquad d=64.
\end{equation}

\subsubsection{多头自注意力编码器}

对每一层 Transformer 编码器，采用标准 post-norm 结构。设第 $l$ 层输入为 $\bm{Z}_{l-1}$，则：
\begin{align}
    \bm{Q}_h &= \bm{Z}_{l-1}\bm{W}_h^Q, \\
    \bm{K}_h &= \bm{Z}_{l-1}\bm{W}_h^K, \\
    \bm{V}_h &= \bm{Z}_{l-1}\bm{W}_h^V,
\end{align}
第 $h$ 个头的注意力为：
\begin{equation}
    \operatorname{head}_h
    =
    \operatorname{softmax}\!\left(
        \frac{\bm{Q}_h \bm{K}_h^\top}{\sqrt{d_h}}
    \right)\bm{V}_h.
\end{equation}

多头拼接后：
\begin{equation}
    \operatorname{MHA}(\bm{Z}_{l-1})
    =
    \operatorname{Concat}(\operatorname{head}_1,\dots,\operatorname{head}_4)\bm{W}^O.
\end{equation}

结合残差与层归一化：
\begin{align}
    \bm{U}_l &= \operatorname{LayerNorm}(\bm{Z}_{l-1} + \operatorname{MHA}(\bm{Z}_{l-1})), \\
    \bm{Z}_l &= \operatorname{LayerNorm}(\bm{U}_l + \operatorname{FFN}(\bm{U}_l)).
\end{align}

代码中共有两层此类编码器，即 $l=1,2$。

\subsubsection{输出投影}

最终通过线性层恢复到变量维度：
\begin{equation}
    \hat{\bm{X}} = \bm{Z}_2 \bm{W}_{\text{out}} + \bm{b}_{\text{out}},
    \qquad \bm{W}_{\text{out}} \in \mathbb{R}^{d \times N}.
\end{equation}

\section{关于 \texorpdfstring{$\bm{X}^{\text{filled}}$}{Xfilled} 的实现细节}

\subsection{代码中的真实写法}

在 \texttt{AGF\_ADNet.forward()} 中，与插补有关的代码逻辑是：
\begin{align}
    \bm{O} &= 1 - \bm{M}, \\
    \bm{O} &= \bm{M}, \\
    \bm{X}^{\text{filled}} &= \bm{X}^{\text{in}} \odot \bm{O} + \bm{X}^{\text{imp}} \odot \bm{M}.
\end{align}

由于第二行直接覆盖了第一行，因此代码真实执行的是：
\begin{equation}
    \bm{X}^{\text{filled}} = \bm{X}^{\text{in}} \odot \bm{M} + \bm{X}^{\text{imp}} \odot \bm{M}.
\end{equation}

\subsection{进一步化简}

又因为进入模型前已经做过
\begin{equation}
    \bm{X}^{\text{in}} = \bm{X} \odot (1-\bm{M}),
\end{equation}
所以有
\begin{equation}
    \bm{X}^{\text{in}} \odot \bm{M} = \bm{0}.
\end{equation}
因此当前代码下 Transformer 实际接收的是：
\begin{equation}
    \bm{X}^{\text{filled}} = \bm{X}^{\text{imp}} \odot \bm{M}.
\end{equation}

\paragraph{含义}
这表示：
\begin{itemize}
    \item 被标记为 missing 的位置进入 Transformer 时保留插补值；
    \item 被标记为 observed 的位置进入 Transformer 时为 0。
\end{itemize}

\paragraph{与常见插补逻辑的差异}
更常见、也更直觉的写法通常是：
\begin{equation}
    \bm{X}^{\text{filled,ideal}}
    =
    \bm{X}^{\text{in}} \odot (1-\bm{M})
    +
    \bm{X}^{\text{imp}} \odot \bm{M},
\end{equation}
即“观测位置保留原值，缺失位置填补预测值”。但本文档在其余部分仍以\textbf{代码真实行为}为准。

\section{损失函数：dual\_domain\_loss}

\subsection{监督目标掩码}

训练时，损失并不是在全部位置上计算，而是只在 \texttt{target\_mask} 指定的位置上监督：
\begin{equation}
    T_{b,s,n} \in \{0,1\}.
\end{equation}

在代码里，$T$ 通常来自“从原本可观测位置中再随机遮掉 $10\%$”这一策略；若某个 batch 恰好一个位置都没抽中，则退化为对全部原始可观测位置监督。

\subsection{时域重构损失}

重构误差项为：
\begin{equation}
    \mathcal{L}_{\text{rec}}
    =
    \frac{
        \sum_{b,s,n}
        \left(
            \hat{X}_{b,s,n} - X_{b,s,n}
        \right)^2
        T_{b,s,n}
    }{
        \max\!\left(\sum_{b,s,n}T_{b,s,n}, 1\right)
    }.
\end{equation}

\subsection{有效样本筛选}

频域损失并非对所有样本都算，而是先基于原始缺失掩码计算每个样本的观测比例：
\begin{equation}
    r_b = \frac{1}{SN}\sum_{s=1}^{S}\sum_{n=1}^{N}(1-M_{b,s,n}).
\end{equation}

仅当
\begin{equation}
    r_b > 0.5
\end{equation}
时，该样本才进入频域损失计算集合：
\begin{equation}
    \mathcal{B}_{\text{valid}} = \{b \mid r_b > 0.5\}.
\end{equation}

\subsection{频域损失}

代码先构造一个频域目标：
\begin{equation}
    \bm{X}^{\text{freq-target}}
    =
    \operatorname{stopgrad}(\hat{\bm{X}})\odot(1-\bm{T})
    +
    \bm{X}\odot\bm{T}.
\end{equation}

这意味着：
\begin{itemize}
    \item 在有监督的位置，用真实值 $\bm{X}$；
    \item 在无监督的位置，用当前预测值的 detached 版本；
    \item 因而频域损失不会把无监督区域的误差梯度反传回去。
\end{itemize}

对有效样本，比较重构与目标的幅值谱：
\begin{align}
    \hat{\bm{F}}_b &= \operatorname{rFFT}(\hat{\bm{X}}_b^\top), \\
    \bm{F}^{\star}_b &= \operatorname{rFFT}((\bm{X}^{\text{freq-target}}_b)^\top), \\
    \mathcal{L}_{\text{freq}}
    &=
    \frac{1}{|\mathcal{B}_{\text{valid}}|}
    \sum_{b \in \mathcal{B}_{\text{valid}}}
    \frac{1}{NF}
    \sum_{n=1}^{N}\sum_{f=1}^{F}
    \left(
        \left|\hat{F}_{b,n,f}\right|
        -
        \left|F^{\star}_{b,n,f}\right|
    \right)^2.
\end{align}

若没有任何有效样本，则代码直接令：
\begin{equation}
    \mathcal{L}_{\text{freq}} = 0.
\end{equation}

\subsection{图熵正则}

尽管函数参数名叫 \texttt{sparsity\_weight}，当前代码使用的并不是 $L_1$ 稀疏正则，而是邻接矩阵的\textbf{行熵}：
\begin{equation}
    \mathcal{L}_{\text{ent}}
    =
    -\frac{1}{N}\sum_{i=1}^{N}\sum_{j=1}^{N}A_{ij}\log(A_{ij}+\varepsilon).
\end{equation}

由于 $\bm{A}$ 已经逐行归一化，这个量越小，说明每一行分布越尖锐，图结构越接近“只关注少数关键邻居”。

\subsection{总损失}

最终损失为：
\begin{equation}
    \mathcal{L}
    =
    \mathcal{L}_{\text{rec}}
    +
    0.1\,\mathcal{L}_{\text{freq}}
    +
    0.1\,\mathcal{L}_{\text{ent}}.
\end{equation}

\section{训练流程的数学化描述}

\subsection{随机再遮挡策略}

设原始缺失掩码为 $\bm{M}$，原始可观测掩码为：
\begin{equation}
    \bm{O} = 1 - \bm{M}.
\end{equation}

代码在可观测位置上额外随机采样一个遮挡指示：
\begin{equation}
    D_{b,s,n} \sim \operatorname{Bernoulli}(0.1),
    \qquad
    \bm{R} = \bm{D} \odot \bm{O}.
\end{equation}

因此真正被随机遮掉的位置是原本可观测且被采样命中的位置。监督掩码定义为：
\begin{equation}
    \bm{T} = \bm{R}.
\end{equation}

若本批次中 $\sum \bm{T}=0$，代码退化为：
\begin{equation}
    \bm{T} = \bm{O}.
\end{equation}

\subsection{构造输入掩码与输入张量}

输入掩码是对原始缺失掩码的拷贝，并把随机遮挡位置也设为缺失：
\begin{equation}
    \bm{M}^{\text{in}} = \bm{M}, \qquad
    M^{\text{in}}_{b,s,n} = 1 \;\text{if}\; T_{b,s,n}=1.
\end{equation}

模型真实接收的输入是：
\begin{equation}
    \bm{X}^{\text{in}} = \bm{X} \odot (1-\bm{M}^{\text{in}}).
\end{equation}

\paragraph{一个重要细节}
\texttt{dual\_domain\_loss()} 接收到的 \texttt{missing\_mask} 仍然是原始 $\bm{M}$，不是 $\bm{M}^{\text{in}}$。也就是说：
\begin{itemize}
    \item 训练输入会看到“原始缺失 + 人工再遮挡”；
    \item 但频域有效样本筛选依据的是“原始缺失比例”。
\end{itemize}

\subsection{优化}

每个 batch 的优化步骤是：
\begin{enumerate}
    \item 前向传播得到 $\hat{\bm{X}}, \bm{A}$；
    \item 计算总损失 $\mathcal{L}$；
    \item 反向传播；
    \item 进行梯度裁剪：
    \begin{equation}
        \|\nabla\|_2 \leftarrow \min(\|\nabla\|_2, 1.0);
    \end{equation}
    \item 由 Adam 更新参数。
\end{enumerate}

每个 epoch 结束后，脚本使用平均训练损失作为 \texttt{ReduceLROnPlateau} 的输入，并保存训练损失最低时的模型参数快照。

\paragraph{实现事实}
这里没有独立验证集；所谓“best model”是按训练集平均损失选出来的。

\section{推理、异常分数与分类评价}

\subsection{异常分数}

推理阶段，模型对每个窗口输出重构结果 $\hat{\bm{X}}$。异常分数定义为原始可观测位置上的平均平方误差：
\begin{equation}
    \operatorname{Score}_b
    =
    \frac{
        \sum_{s=1}^{S}\sum_{n=1}^{N}
        (\hat{X}_{b,s,n} - X_{b,s,n})^2 (1-M_{b,s,n})
    }{
        \sum_{s=1}^{S}\sum_{n=1}^{N}(1-M_{b,s,n}) + 10^{-8}
    }.
\end{equation}

这里分母的 \texttt{clamp\_min(1e-8)} 用来避免极端情况下除零。

\subsection{阈值估计}

训练完成后，脚本先在\textbf{训练窗口}上计算一遍异常分数：
\begin{equation}
    \{\operatorname{Score}^{\text{train}}_1, \dots, \operatorname{Score}^{\text{train}}_K\}.
\end{equation}

阈值定义为：
\begin{equation}
    \tau
    =
    \operatorname{mean}(\operatorname{Score}^{\text{train}})
    +
    \operatorname{std}(\operatorname{Score}^{\text{train}}).
\end{equation}

\paragraph{实现事实}
这不是三倍标准差规则，而是一倍标准差规则，即 \texttt{mean + std}。

\subsection{测试标签构造与评价}

\texttt{build\_test\_labels(num\_scores)} 直接把测试集后半段标为异常：
\begin{equation}
    y_i =
    \begin{cases}
        0, & i < \lfloor K/2 \rfloor \\
        1, & i \ge \lfloor K/2 \rfloor
    \end{cases}
\end{equation}

于是预测标签为：
\begin{equation}
    \hat{y}_i = \mathbb{I}[\operatorname{Score}_i > \tau].
\end{equation}

随后计算准确率、精确率、召回率和 F1：
\begin{align}
    \operatorname{Accuracy} &= \frac{TP + TN}{TP + TN + FP + FN}, \\
    \operatorname{Precision} &= \frac{TP}{TP + FP}, \\
    \operatorname{Recall} &= \frac{TP}{TP + FN}, \\
    F_1 &= \frac{2 \cdot \operatorname{Precision} \cdot \operatorname{Recall}}{\operatorname{Precision} + \operatorname{Recall}}.
\end{align}

\paragraph{实现事实}
这种标签构造方式隐含了一个强假设：测试窗口在时间顺序上“前半正常、后半异常”。

\section{逐层张量数据流表}

\subsection{模型前向阶段}

\begin{longtable}{>{\raggedright\arraybackslash}p{0.22\linewidth} >{\raggedright\arraybackslash}p{0.28\linewidth} >{\raggedright\arraybackslash}p{0.28\linewidth}}
\toprule
阶段 & 张量表达式 & 形状 \\
\midrule
\endhead
输入窗口 & $\bm{X}$ & $B \times S \times N$ \\
输入掩码 & $\bm{M}^{\text{in}}$ & $B \times S \times N$ \\
缺失置零 & $\bm{X}^{\text{in}} = \bm{X}\odot(1-\bm{M}^{\text{in}})$ & $B \times S \times N$ \\
图学习输出 & $\bm{A}$ & $N \times N$ \\
GCN 输入扩维 & $\bm{X}^{\text{in}}.\text{unsqueeze}(-1)$ & $B \times S \times N \times 1$ \\
GCN 输出 & $\bm{H}^{\text{gcn}}$ & $B \times S \times N$ \\
TCN 转置输入 & $\operatorname{permute}(\bm{X}^{\text{in}})$ & $B \times N \times S$ \\
三尺度堆叠 & $\bm{Y}^{\text{stack}}$ & $B \times N \times S \times 3$ \\
TCN 输出 & $\bm{H}^{\text{tcn}}$ & $B \times S \times N$ \\
时域融合 & $\bm{H}^{\text{time}}$ & $B \times S \times N$ \\
频域 FFT 输入 & $\operatorname{permute}(\bm{X}^{\text{in}})$ & $B \times N \times S$ \\
频谱张量 & $\bm{X}_f$ & $B \times N \times F$（复数） \\
频域输出 & $\bm{H}^{\text{freq}}$ & $B \times S \times N$ \\
门控拼接张量 & $\bm{C}$ & $B \times 2N \times S$ \\
门控图 & $\bm{Z}_g$ & $B \times S \times N$ \\
插补结果 & $\bm{X}^{\text{imp}}$ & $B \times S \times N$ \\
Transformer 输入 & $\bm{X}^{\text{filled}}$ & $B \times S \times N$ \\
输入投影后 & $\bm{Z}_0$ & $B \times S \times 64$ \\
Transformer 输出 & $\bm{Z}_2$ & $B \times S \times 64$ \\
最终重构 & $\hat{\bm{X}}$ & $B \times S \times N$ \\
\bottomrule
\end{longtable}

\subsection{训练外层流程}

\begin{longtable}{>{\raggedright\arraybackslash}p{0.18\linewidth} >{\raggedright\arraybackslash}p{0.70\linewidth}}
\toprule
阶段 & 数据流说明 \\
\midrule
\endhead
CSV 读取 & 读取 \texttt{./data/train/train\_*.csv} 与 \texttt{./data/test/test\_*.csv}，生成原始数组与缺失掩码。 \\
标准化 & 在训练集上拟合 \texttt{StandardScaler}，对训练与测试都做同一组变换。 \\
窗口化 & 构造长度 60 的滑动窗口，序列开头不足长度时用首行复制填充。 \\
训练输入生成 & 在原始可观测位置中随机遮掉 10\%，形成 \texttt{m\_input} 与 \texttt{x\_input}。 \\
模型前向 & \texttt{AGF\_ADNet(x\_input, m\_input)} 输出重构序列、图结构与插补表示。 \\
损失计算 & 用原始窗口 \texttt{x}、原始缺失掩码 \texttt{m}、目标掩码 \texttt{target\_mask} 计算双域损失。 \\
最优模型保存 & 按平均训练损失选取最佳参数快照。 \\
训练阈值估计 & 对训练窗口跑一遍异常分数，阈值取 \texttt{mean + std}。 \\
测试集评价 & 对测试窗口计算分数，再和阈值比较，输出准确率、精确率、召回率、F1。 \\
\bottomrule
\end{longtable}

\section{实现层面的关键观察}

\begin{enumerate}
    \item \textbf{掩码语义是反直觉的。} 在整个脚本里，$1$ 表示 missing，$0$ 表示 observed。
    \item \textbf{标准化发生在零填充之后。} 缺失值在标准化前先被置为 0，再参与均值方差计算。
    \item \textbf{频域分支真正使用的是实部/虚部。} 幅值与相位虽然被计算，但没有直接参与增强或损失。
    \item \textbf{图正则不是 $L_1$，而是熵。} 这与很多图学习论文中的“稀疏正则”写法不同。
    \item \textbf{阈值是 $1\sigma$ 规则。} 代码实现为训练分数均值加一个标准差，而非三倍标准差。
    \item \textbf{测试标签生成依赖顺序假设。} 后半段强制视为异常，是一种数据集特定约定，不是通用异常标注方式。
    \item \textbf{当前 Transformer 输入只保留掩码位置的插补值。} 这是由 \texttt{observed\_mask} 覆盖造成的真实行为。
    \item \textbf{最佳模型按训练损失选取。} 没有单独验证集，也没有早停机制。
\end{enumerate}

\section{总结}

\texttt{mra.py} 实现的是一个以\textbf{重构误差驱动异常检测}为核心思想的模型，主要由以下几层机制叠加而成：
\begin{enumerate}
    \item 用图学习模块自适应建模变量间依赖；
    \item 用 GCN 与多尺度 TCN 提取时域结构；
    \item 用 FFT、频率注意力和频域残差增强提取频谱结构；
    \item 用门控机制在时域与频域之间做动态融合；
    \item 用 Transformer 做全局时序重构；
    \item 用时域误差、频域谱误差和图熵约束联合训练；
    \item 用训练集分数的均值加标准差作为异常阈值。
\end{enumerate}

如果把它抽象成一句话，那么这份代码的数学本质可以概括为：
\begin{quote}
    先把多变量时间序列在\textbf{图结构、时域局部模式、频域全局模式}三个视角下进行表征，再通过融合与 Transformer 重构建立“正常行为模型”，最后用\textbf{掩码加权的重构误差}判断异常。
\end{quote}

但在实际部署或复现实验时，仍应特别注意本文档指出的实现细节，尤其是掩码语义、阈值公式以及 \texttt{X\_filled} 的当前代码行为。

\end{document}
